\section{Results I: At-Site Failures Are Repaired Locally}

\subsection{Volume-Control Ablation}

The first experiment asks whether supervision supplies enough intermediate
volume control. It compares a baseline MLP, NLS alone, partial volume control,
full volume control, and volume control without NLS on MNIST with $K=25$.

\begin{table}[h]
\centering
\begin{tabular}{lrrrrr}
\toprule
Config & Dead & Min Var & Redundancy & Entropy & Accuracy \\
\midrule
Baseline & 1.40 & 0.000 & 20.5 & 3.08 & 96.09 \\
NLS only & 1.20 & 0.264 & 24.9 & 2.93 & 96.14 \\
NLS + Var & 0.00 & 11.74 & 187.0 & 2.09 & 95.75 \\
NLS + Var + Decorr & 0.00 & 7.84 & 13.9 & 2.46 & 96.34 \\
Var + Decorr only & 0.10 & 3.55 & 13.3 & 2.65 & 96.40 \\
\bottomrule
\end{tabular}
\caption{Volume-control ablation summary. Accuracy is reported in percent.}
\label{tab:volume-ablation}
\end{table}

Supervised gradients do not provide intermediate volume control: the baseline
has dead ReLU units in every seed. NLS without volume control is insufficient.
Variance alone is pathological, producing alive but highly redundant units.
Full volume control repairs the site. Volume control without NLS is useful
regularization, but it does not produce the same calibrated EM structure.

\subsection{Capacity Sweep}

The capacity sweep compares the baseline with NLS plus full volume control for
$K \in \{16,25,36,49,64\}$. The baseline has dead units at every width, with a
stable dead fraction around five to seven percent. Volume control eliminates
dead units at every width. Redundancy grows with width, and volume control
reduces the slope and tightens seed variability.

The accuracy effect is capacity- and $\lambda$-dependent. This paper therefore
uses the capacity experiment as evidence for structural health rather than as a
benchmark claim.
